\bta{2010}


\begin{enumerate}
	%\renewcommand{\labelenumi}{\arabic{enumi}.}
	% A(\Alph) a(\alph) I(\Roman) i(\roman) 1(\arabic)
	%设定全局标号series=example	%引用全局变量resume=example
	%[topsep=-0.3em,parsep=-0.3em,itemsep=-0.3em,partopsep=-0.3em]
	%可使用leftmargin调整列表环境左边的空白长度 [leftmargin=0em]
	\item


\begin{enumerate}
	%\renewcommand{\labelenumi}{\arabic{enumi}.}
	% A(\Alph) a(\alph) I(\Roman) i(\roman) 1(\arabic)
	%设定全局标号series=example	%引用全局变量resume=example
	%[topsep=-0.3em,parsep=-0.3em,itemsep=-0.3em,partopsep=-0.3em]
	%可使用leftmargin调整列表环境左边的空白长度 [leftmargin=0em]
	\item
设$\hat{A},\hat{B}$与Pauli算符对易，证明：
$(\vec{\sigma}\cdot\vec{A})(\vec{\sigma}\cdot\vec{B})=\vec{A}\cdot\vec{B}+i\vec{\sigma}\cdot(\vec{A}\times\vec{B})$

\item 
试将$(\hat{I}+\hat{\sigma}_x+i\hat{\sigma}_y)^{\frac{1}{2}}$表示成$\hat{I},\hat{\sigma}_x,\hat{\sigma}_y,\hat{\sigma_z}$的线性叠加，$\hat{I}$为单位算符。

\end{enumerate}
\banswer{
\begin{enumerate}
	%\renewcommand{\labelenumi}{\arabic{enumi}.}
	% A(\Alph) a(\alph) I(\Roman) i(\roman) 1(\arabic)
	%设定全局标号series=example	%引用全局变量resume=example
	%[topsep=-0.3em,parsep=-0.3em,itemsep=-0.3em,partopsep=-0.3em]
	%可使用leftmargin调整列表环境左边的空白长度 [leftmargin=0em]
		\item
		Pauli 矩阵满足
		\[
			\sigma_i\sigma_j=\delta_{ij}I+i\sum_k\varepsilon_{ijk}\sigma_k.
		\]
		又因 $\vec A,\vec B$ 的各分量与 Pauli 矩阵对易，
		\begin{align*}
			(\vec\sigma\cdot\vec A)(\vec\sigma\cdot\vec B)
			&=\sum_{ij}A_iB_j\sigma_i\sigma_j\\
			&=\sum_iA_iB_i
			+i\sum_{ijk}\varepsilon_{ijk}A_iB_j\sigma_k\\
			&=\boxed{\vec A\cdot\vec B
			+i\vec\sigma\cdot(\vec A\times\vec B)}.
		\end{align*}
		\item
		令 $N=\sigma_x+i\sigma_y$。由 Pauli 代数可得
		$N^2=0$，所以
		\[
			\left(I+\frac12N\right)^2=I+N.
		\]
		因此一个平方根为
		$\sqrt{I+\sigma_{x}+i \sigma_{y}}=I+\frac{1}{2}\left(\sigma_{x}+i \sigma_{y}\right)$
		或者
		$\sqrt{I+\sigma_{x}+i \sigma_{y}}=-I-\frac{1}{2}\left(\sigma_{x}+i \sigma_{y}\right)$
	
\end{enumerate}
\begin{note}
	表示方法不唯一，这是由于根号的多值性决定的
\end{note}

}




\item 
设一维谐振子的初态为$\psi(x,0)=\cos\frac{\theta}{2}\varphi_0(x)+\sin\frac{\theta}{2}\varphi_1(x)$，
即基态与第一激发态的叠加，其中$\theta$为实参数：
\begin{enumerate}
	%\renewcommand{\labelenumi}{\arabic{enumi}.}
	% A(\Alph) a(\alph) I(\Roman) i(\roman) 1(\arabic)
	%设定全局标号series=example	%引用全局变量resume=example
	%[topsep=-0.3em,parsep=-0.3em,itemsep=-0.3em,partopsep=-0.3em]
	%可使用leftmargin调整列表环境左边的空白长度 [leftmargin=0em]
	\item
求$t$时刻的波函数$\psi(x,t)$；

\item 
求$t$时刻粒子处于基态及第一激发态的概率；

\item 
求$t$时刻粒子的势能算符$\hat{V}=\frac{m\omega^2}{2}x^2$的平均值；

\item 
求演化成$-\psi(x,0)$所需的最短时间$t_{min}$。
\end{enumerate}
\banswer{
令 $c=\cos(\theta/2)$、$s=\sin(\theta/2)$，谐振子能量为
$E_n=(n+1/2)\hbar\omega$。
\begin{enumerate}
\item 各定态分别积累动力学相位，故位置表象的波函数为
\[
 \boxed{\psi(x,t)=c\varphi_0(x)e^{-i\omega t/2}
                    +s\varphi_1(x)e^{-3i\omega t/2}}.
\]
\item 利用本征态的正交归一性，投影得
$\boxed{P_0=|c|^2,\ P_1=|s|^2}$；相位不影响能量概率。
\item 写 $x=x_0(a+a^\dagger)$，$x_0=\sqrt{\hbar/(2m\omega)}$。
于是
\[
 x^2=x_0^2(a^2+a^{\dagger2}+aa^\dagger+a^\dagger a),\qquad
 \langle n|x^2|n\rangle=x_0^2(2n+1).
\]
$x^2$ 只能使量子数改变 $0,\pm2$，故 $\langle0|x^2|1\rangle=0$，
也可由宇称判断交叉项为零。因此
\[
 \boxed{\langle V\rangle_t=\frac{\hbar\omega}{4}c^2
                             +\frac{3\hbar\omega}{4}s^2}.
\]
\item 若 $c,s$ 均非零，要求 $\psi(x,t)=-\psi(x,0)$，须同时满足
\[
 e^{-i\omega t/2}=-1,\qquad e^{-3i\omega t/2}=-1.
\]
第一式要求 $\omega t/2=(2k+1)\pi$，此时第二式自动满足，故
$\boxed{t_{\min}=2\pi/\omega}$。
若 $s=0$，结果仍为 $2\pi/\omega$；若 $c=0$，初态只有 $n=1$，则最短时间改为
$\boxed{2\pi/(3\omega)}$。这里要求波函数严格变号；整体负号并不改变物理态。
\end{enumerate}
}



\item 
设基态氢原子处于弱电场中，微扰哈密顿量为：
$$
\hat{H} ^{\prime} =\begin{cases}0,&t\le0\\
\lambda ze^{-\frac{t}{T}},&t>0
\end{cases}
$$
其中$  \lambda,T $为常数。
\begin{enumerate}
	%\renewcommand{\labelenumi}{\arabic{enumi}.}
	% A(\Alph) a(\alph) I(\Roman) i(\roman) 1(\arabic)
	%设定全局标号series=example	%引用全局变量resume=example
	%[topsep=-0.3em,parsep=-0.3em,itemsep=-0.3em,partopsep=-0.3em]
	%可使用leftmargin调整列表环境左边的空白长度 [leftmargin=0em]
	\item
求很长时间后$(t\gg T)$电子从基态跃迁到第一激发态的概率。
基态和激发态波函数为：
$$
\begin{aligned}
	&\psi_{100}(\vec{r})=R_{10}(r)Y_{00}(\theta,\varphi)=\frac{2}{a^{\frac{3}{2}}}e^{-\frac{r}{a}} \cdot \frac{1}{\sqrt{4\pi}} 
	\\ 
	&\psi_{210}(\vec{r})=R_{21}(r)Y_{10}(\theta,\varphi)=\frac{1}{(2a)^{\frac{3}{2}}}\frac{r}{\sqrt3a}e^{-\frac{r}{2a}} \cdot \sqrt{\frac{3}{4\pi}}\cos\theta 
\end{aligned}
$$
其中$  a  $为玻尔半径。
\item 
基态电子跃迁到下列哪个激发态的概率等于零？简述理由。
\fourchoices
{$\psi_{200}$}
{$\psi_{211}$}
{$\psi_{21,-1}$}
{$\psi_{210}$}


	
\end{enumerate}

\banswer{
令 $\omega_{21}=(E_2-E_1)/\hbar$。微扰的空间部分为 $\lambda r\cos\theta$。
\begin{enumerate}
\item 由 $\cos\theta\,Y_{00}=Y_{10}/\sqrt3$，角积分为
$\int Y_{10}^*\cos\theta Y_{00}\,d\Omega=1/\sqrt3$。
径向积分含体积元 $r^2dr$ 和算符中的一个 $r$，故
\begin{align*}
 z_{21}&=\langle210|z|100\rangle
 =\frac1{\sqrt3}\int_0^\infty R_{21}(r)R_{10}(r)r^3dr\\
 &=\frac1{3\sqrt2a^4}\int_0^\infty r^4e^{-3r/(2a)}dr
 =\frac{4!}{3\sqrt2a^4}\left(\frac{2a}{3}\right)^5
 =\frac{128\sqrt2}{243}a.
\end{align*}
因此 $H'_{21}(t)=\lambda z_{21}e^{-t/T}$，其中 $\lambda$ 不能遗漏。
一级含时微扰论给出
\begin{align*}
 c_{210}^{(1)}(\infty)
 &=-\frac{i\lambda z_{21}}{\hbar}\int_0^\infty e^{-(1/T-i\omega_{21})t}dt\\
 &=-\frac{i\lambda z_{21}T}{\hbar(1-i\omega_{21}T)}.
\end{align*}
故 $t\gg T$ 时
\[
 \boxed{P_{100\to210}=\frac{\lambda^2|z_{21}|^2T^2}
 {\hbar^2(1+\omega_{21}^2T^2)}
 =\frac{2^{15}}{3^{10}}\frac{\lambda^2a^2T^2}
 {\hbar^2(1+\omega_{21}^2T^2)}}.
\]
\item 角向积分要求 $\Delta l=\pm1$、$\Delta m=0$。
从 $l=m=0$ 出发，$n=2$ 子空间只有 $|210\rangle$ 被耦合。
$|200\rangle$ 的宇称不合，$|211\rangle$、$|21,-1\rangle$ 的磁量子数不合，
所以这三个选项的一级跃迁振幅均为零。第一激发能级的总跃迁概率即上一问的概率。
这些选择定则针对当前的一阶电偶极跃迁；高阶过程不能由此一概排除。
\end{enumerate}
}


\newpage

\item 
两个质量为$m$的粒子处于一个边长为$a>b>c$的不可穿透的长盒子中，求下列条件下，该体系能量最低态的波函数（只写出空间部分）及对应能量。
\begin{enumerate}
	%\renewcommand{\labelenumi}{\arabic{enumi}.}
	% A(\Alph) a(\alph) I(\Roman) i(\roman) 1(\arabic)
	%设定全局标号series=example	%引用全局变量resume=example
	%[topsep=-0.3em,parsep=-0.3em,itemsep=-0.3em,partopsep=-0.3em]
	%可使用leftmargin调整列表环境左边的空白长度 [leftmargin=0em]
	\item
非全同粒子。
\item 
零自旋全同粒子。
\item 
自旋为$\frac{1}{2}$的全同粒子。
\end{enumerate}
\banswer{
一个粒子处于长盒子中的波函数及相应能量为
$$
\phi(x, y, z)=\frac{2 \sqrt{2}}{\sqrt{a b c}} \sin \frac{n_{1} \pi x}{a} \sin \frac{n_{2} \pi y}{b} \sin \frac{n_{3} \pi z}{c} \quad e_{n}=\frac{\pi^{2} \hbar^{2}}{2 m}\left(\frac{n_{1}^{2}}{a^{2}}+\frac{n_{2}^{2}}{b^{2}}+\frac{n_{3}^{2}}{c^{2}}\right)
$$
\begin{enumerate}
	%\renewcommand{\labelenumi}{\arabic{enumi}.}
	% A(\Alph) a(\alph) I(\Roman) i(\roman) 1(\arabic)
	%设定全局标号series=example	%引用全局变量resume=example
	%[topsep=-0.3em,parsep=-0.3em,itemsep=-0.3em,partopsep=-0.3em]
	%可使用leftmargin调整列表环境左边的空白长度 [leftmargin=0em]
	\item
	非全同粒子基态：
	$$\psi_{1}\left(r_{1}, r_{2}\right)=\phi_{1}\left(r_{1}\right) \phi_{1}\left(r_{2}\right)=\frac{8}{a b c} \sin \frac{\pi x_{1}}{a} \sin \frac{\pi y_{1}}{b} \sin \frac{\pi z_{1}}{c} \sin \frac{\pi x_{2}}{a} \sin \frac{\pi y_{2}}{b} \sin \frac{\pi z_{2}}{c}$$
	$$
	E=e_{1}+e_{1}=\frac{\pi^{2} \hbar^{2}}{m}\left(\frac{1}{a^{2}}+\frac{1}{b^{2}}+\frac{1}{c^{2}}\right)
	$$
	\item 
	零自旋全同粒子基态，零自旋粒子可以处于同一态，与非全同粒子基态相同：
	$$
	\psi_{1}\left(r_{1}, r_{2}\right)=\frac{1}{2}\left[\phi_{1}\left(r_{1}\right) \phi_{1}\left(r_{2}\right)+\phi_{1}\left(r_{2}\right) \phi_{1}\left(r_{1}\right)\right]=\phi_{1}\left(r_{1}\right) \phi_{1}\left(r_{2}\right)
	$$
	\item 
		自旋为 $\frac{1}{2}$ 的全同费米子若取自旋单态，则自旋部分反对称，空间部分可取对称基态。因此体系的真正基态仍为
		$$
		\psi_{\mathrm{space}}(r_1,r_2)=\phi_1(r_1)\phi_1(r_2),
		\qquad
		E=2e_1=\frac{\pi^{2} \hbar^{2}}{m}
		\left(\frac{1}{a^{2}}+\frac{1}{b^{2}}+\frac{1}{c^{2}}\right).
		$$
		若另行限定两电子处于自旋三重态，则空间部分必须反对称。此时由于 $a>b>c$，最低的第二个单粒子轨道为
		$$
		\phi_{2}(r)=\frac{2 \sqrt{2}}{\sqrt{a b c}}
		\sin \frac{2 \pi x}{a}\sin \frac{\pi y}{b}\sin \frac{\pi z}{c},
		$$
		并有
		$$
		\psi_{\mathrm{space}}=\frac{1}{\sqrt{2}}
		\left[\phi_{1}(r_{1}) \phi_{2}(r_{2})
		-\phi_{2}(r_{1}) \phi_{1}(r_{2})\right],
		\quad
		E=\frac{\pi^{2} \hbar^{2}}{m}
		\left(\frac{5}{2a^{2}}+\frac{1}{b^{2}}+\frac{1}{c^{2}}\right).
$$
\end{enumerate}
\begin{note}
此题要求只写出空间部分，一般情况下还有自旋部分。当盒子是方盒子即 $a=b=c$ 时，第二激发态$\phi_{2}$ 存在简并，三重，对应的波函数有$ 3 $个。另外波函数的形式随着坐标原点的选择不同而不同（也可将坐标原点选在中心）。
\end{note}
\begin{note}
此题第三问在考虑粒子自旋的情况下基态是与前两问相同的，只有在限定为三重态的情况下才会改变基态，此题正确题目及结果应参照陈书7.7
\end{note}

}



\item 
粒子在宽度为$ 2a $的一维无限深势阱中运动
\[ 
V(x)=
\left\{
\begin{array}{lr}
	0,&0<x<2a\\
	\infty, & x\leq 0,x\geq2a
\end{array}
\right.
 \]
设该体系受到$\hat{H} ^{\prime} =\lambda\delta(x-a)$的微扰作用。
\begin{enumerate}
	%\renewcommand{\labelenumi}{\arabic{enumi}.}
	% A(\Alph) a(\alph) I(\Roman) i(\roman) 1(\arabic)
	%设定全局标号series=example	%引用全局变量resume=example
	%[topsep=-0.3em,parsep=-0.3em,itemsep=-0.3em,partopsep=-0.3em]
	%可使用leftmargin调整列表环境左边的空白长度 [leftmargin=0em]
	\item
利用微扰理论，求第$n$能级精确到二级的近似表达式。

\item 
指出所得结果的适用条件。
\end{enumerate}

\banswer{
无扰本征态和能量为
\[
 \phi_n(x)=\frac1{\sqrt a}\sin\frac{n\pi x}{2a},\qquad
 E_n^{(0)}=n^2\epsilon,\quad \epsilon=\frac{\pi^2\hbar^2}{8ma^2}.
\]
由 $\delta$ 函数的取样性质，先求所有微扰矩阵元：
\[
 H'_{kn}=\lambda\phi_k(a)\phi_n(a)
        =\frac{\lambda}{a}\sin\frac{k\pi}{2}\sin\frac{n\pi}{2}.
\]
若 $k,n$ 中任一个为偶数，矩阵元为零；若二者均为奇数，矩阵元为
$\pm\lambda/a$，符号由两个正弦的乘积决定，不能统一写为正号。
\begin{enumerate}
\item 非简并微扰公式给出
\[
 E_n^{(1)}=H'_{nn},\qquad
 E_n^{(2)}=\sum_{k\ne n}\frac{|H'_{kn}|^2}{\epsilon(n^2-k^2)}.
\]
对偶数 $n$，$\phi_n(a)=0$，微扰作用于该态恒为零，因此能量严格不变：
\[
 \boxed{E_n=n^2\epsilon\quad(n\text{ 为偶数})}.
\]
对奇数 $n$，$E_n^{(1)}=\lambda/a$，二阶只需对奇数 $k$ 求和。
利用正切函数的部分分式展开，
\[
 \sum_{k=1,3,5,\ldots}\frac1{z^2-k^2}
 =-\frac{\pi}{4z}\tan\frac{\pi z}{2}.
\]
令 $z=n+\delta$，减掉 $k=n$ 项再取极限，
\[
 \sum_{\substack{k=1,3,5,\ldots\\k\ne n}}\frac1{n^2-k^2}
 =\lim_{\delta\to0}\left[-\frac{\pi}{4(n+\delta)}
 \tan\frac{\pi(n+\delta)}2-\frac1{(n+\delta)^2-n^2}\right]
 =-\frac1{4n^2}.
\]
这里两项的 $1/\delta$ 奇性相消。代回即得
\[
 \boxed{E_n=\frac{n^2\pi^2\hbar^2}{8ma^2}+\frac{\lambda}{a}
       -\frac{2m\lambda^2}{\pi^2\hbar^2n^2}+O(\lambda^3)
       \quad(n\text{ 为奇数})}.
\]
\item 要求对所有重要的奇数 $k\ne n$，
\[
 \left|\frac{\lambda/a}{E_n^{(0)}-E_k^{(0)}}\right|\ll1.
\]
对低能级，一个充分的尺度条件为 $m|\lambda|a/\hbar^2\ll1$。
偶数态因在微扰位置有节点，其精确不变性不依赖上述弱耦合近似。
\end{enumerate}
}



\end{enumerate}
